\chapter{Límites y continuidad}\label{ch:b1:continuity}

El teorema del valor intermedio y el de los valores extremos se
usaron en el volumen anterior sobre la base de la evidencia visual.
Con las sucesiones (\cref{ch:b1:seq}) y la topología de $\R$
(\cref{ch:b1:topology}) ya en la mano, este capítulo los demuestra ---
y completa la teoría con el teorema de la \omterm{def:b1:logic:inj}{biyección} monótona (que
legitima $\arcsin$, $\operatorname{arcosh}$ y sus parientes
del~\cref{ch:b1:functions}) y con el teorema de Heine sobre la
\omterm{def:b1:continuity:uniform}{continuidad uniforme}.

En todo el capítulo, $I$ es un \omterm{prop:b1:reals:intervals}{intervalo} y $f \colon I \to \R$; escribir
«$x_0 \in \overline I$» permite límites en los extremos.

\section{Límites de funciones}

\begin{definition}[Límite en un punto]\label{def:b1:continuity:limit}
Sean $x_0 \in \overline{I}$ y $\ell \in \R$. Se dice que
$f(x) \to \ell$ cuando $x \to x_0$ si\index{límite!de una función}
\[
\forall \varepsilon > 0,\ \exists \delta > 0,\ \forall x \in I,
\qquad \abs{x - x_0} \leq \delta \implies \abs{f(x) - \ell} \leq
\varepsilon .
\]
Los límites en $\pm\infty$ y los límites infinitos se definen con el
mismo patrón ($\abs{x - x_0} \leq \delta$ se convierte en $x \geq M$;
$\abs{f(x) - \ell} \leq \varepsilon$, en $f(x) \geq M'$). Los límites
laterales restringen $x$ a $x > x_0$ (se escribe $x \to x_0^+$) o a
$x < x_0$. El límite es único cuando existe (misma demostración que
en la~\cref{prop:b1:seq:first}).
\end{definition}

\begin{example}[Un límite por emparedado]\label{ex:b1:continuity:floorlimit}
Calcúlese $\lim_{x \to 0} x\,\bigl\lfloor \frac1x \bigr\rfloor$. El
encuadre de la \omterm{thm:b1:reals:floor}{parte entera} $\frac1x - 1 < \lfloor \frac1x \rfloor
\leq \frac1x$ da, multiplicando por $x$ (¡ojo al signo!):
\[
1 - x < x\Bigl\lfloor \frac1x \Bigr\rfloor \leq 1 \quad (x > 0),
\qquad
1 \leq x\Bigl\lfloor \frac1x \Bigr\rfloor < 1 - x \quad (x < 0),
\]
y los dos emparedados laterales se cierran sobre $1$: el límite es
$1$. Obsérvese lo ocurrido: $\lfloor \frac1x\rfloor$ por sí solo da
saltos salvajes cerca de $0$, pero el factor $x$ domestica cada salto
($x$ por un salto unidad es pequeño) y solo sobrevive el encuadre. La
idea de cierre: los límites de productos de un factor pequeño por un
factor de oscilación acotada son problemas de emparedado, nunca
problemas del teorema de las operaciones --- este último necesita que
converjan \emph{los dos} factores.
\end{example}

\begin{theorem}[Caracterización secuencial]\label{thm:b1:continuity:seqchar}
$f(x) \to \ell$ cuando $x \to x_0$ si y solo si, para \emph{toda}
sucesión $(u_n)$ de puntos de $I$ con $u_n \to x_0$, se tiene
$f(u_n) \to \ell$.
\end{theorem}

\begin{proof}
($\Rightarrow$) Sean $u_n \to x_0$ y $\varepsilon > 0$. Tómese
$\delta$ de la definición y después $N$ con
$\abs{u_n - x_0} \leq \delta$ para $n \geq N$: más allá de $N$,
$\abs{f(u_n) - \ell} \leq \varepsilon$.

($\Leftarrow$) Por contrarrecíproco. Si $f \not\to \ell$, algún
$\varepsilon_0 > 0$ derrota a todo $\delta$; eligiendo
$\delta = \frac{1}{n+1}$ se obtienen $u_n \in I$ con
$\abs{u_n - x_0} \leq \frac{1}{n+1}$ y
$\abs{f(u_n) - \ell} > \varepsilon_0$. Entonces $u_n \to x_0$ pero
$f(u_n) \not\to \ell$.
\end{proof}

\begin{corollary}[Operaciones, composición, orden]\label{cor:b1:continuity:operations}
Las sumas, los productos y los cocientes (con límite no nulo en el
denominador) de límites se comportan como en las sucesiones; si
$f \to \ell$ en $x_0$ y $g \to m$ en $\ell$, y $g$ está definida en
torno a $\ell$ con $g(\ell) = m$, o bien $f \neq \ell$ cerca de
$x_0$, entonces $g \circ f \to m$ en $x_0$; los límites conservan las
desigualdades amplias y vale el teorema del emparedado.
\end{corollary}

\begin{proof}
Cada \omterm{def:b1:logic:statement}{enunciado} se transfiere por
la~\cref{thm:b1:continuity:seqchar} a su análogo para sucesiones
(\cref{thm:b1:seq:operations,thm:b1:seq:order}). Para la composición
merece la pena escribir el encadenamiento directo una vez: dado
$\varepsilon > 0$, el límite de $g$ en $\ell$ proporciona
$\eta > 0$ con
\[
\abs{y - \ell} \leq \eta \implies \abs{g(y) - m} \leq
\varepsilon \quad (y \text{ en el dominio de } g),
\]
donde el caso $y = \ell$ queda cubierto porque $g(\ell) = m$; después
el límite de $f$ en $x_0$ proporciona $\delta > 0$ con
$\abs{x - x_0} \leq \delta \implies \abs{f(x) - \ell} \leq \eta$; y
encadenando, $\abs{x - x_0} \leq \delta$ da
$\abs{g(f(x)) - m} \leq \varepsilon$. Con la hipótesis alternativa
($f \neq \ell$ cerca de $x_0$), el valor $y = \ell$ nunca se le pasa
a $g$ y lo que valga ahí es irrelevante.
\end{proof}

\begin{example}[Por qué existe la salvedad de la composición]\label{ex:b1:continuity:compositiontrap}
Sea $g(y) = 0$ para $y \neq 0$ y $g(0) = 1$, y sea $f$ idénticamente
$0$. Entonces $f(x) \to 0$ cuando $x \to 0$, y $g(y) \to 0$ cuando
$y \to 0$; y, sin embargo, $g(f(x)) = g(0) = 1$ para todo $x$, luego
$g \circ f \to 1 \neq 0$. La función \omterm{def:b1:topology:closure}{interior} se sienta
\emph{exactamente sobre} el valor prohibido $\ell = 0$ para siempre,
y el límite de $g$ en $0$ ignora lo que $g$ hace en $0$. La salvedad
del~\cref{cor:b1:continuity:operations} ---o bien $g(\ell) = m$ (es
decir, $g$ \omterm{def:b1:continuity:continuous}{continua} en $\ell$), o bien $f \neq \ell$ cerca de
$x_0$--- es justo lo que excluye este caso. La idea de cierre: en la
práctica se componen funciones \emph{\omterm{def:b1:continuity:continuous}{continuas}} y la salvedad sale
gratis; solo muerde cuando los límites se toman a lo largo de
\omterm{def:b1:topology:open}{entornos} perforados, y por eso la definición de $\lim_{x \to x_0}$
usada en este libro incluye el punto cuando está en el dominio.
\end{example}

\section{Continuidad}

\begin{definition}\label{def:b1:continuity:continuous}
$f$ es \emph{continua en $x_0 \in I$}\index{continuidad} cuando
$f(x) \to f(x_0)$ al tender $x \to x_0$; y \emph{continua en $I$}
cuando lo es en cada punto. Por la~\cref{thm:b1:continuity:seqchar},
$f$ es continua en $x_0$ si y solo si $f(u_n) \to f(x_0)$ para toda
sucesión $u_n \to x_0$ de $I$.
\end{definition}

\begin{example}[Una taxonomía de las discontinuidades]\label{ex:b1:continuity:taxonomy}
Tres maneras de fallar en un punto, en orden creciente de gravedad.
\emph{Evitable}: $f(x) = \frac{\sin x}{x}$ en $\R^*$ tiene límite $1$
en $0$; definir $f(0) = 1$ lo repara --- la discontinuidad era un
agujero, no un rasgo. \emph{De salto}: $\lfloor x \rfloor$ en un
entero tiene límites laterales distintos ($n - 1$ y $n$); ninguna
elección de valor puede reconciliarlos, pero los dos semilímites
existen. \emph{Esencial}: $\sin\frac1x$ en $0$ no tiene ningún límite
lateral (\cref{exo:b1:continuity:1}) --- oscilación sin
estabilización. Las funciones monótonas solo pueden producir la
variedad intermedia (sus límites laterales existen siempre, por ser
\omterm{def:b1:reals:bounds}{supremos} e \omterm{def:b1:reals:bounds}{ínfimos}), y por eso sus \omterm{def:b1:logic:sets}{conjuntos} de discontinuidad son
a lo sumo numerables --- un racional por salto. Las derivadas, por el
teorema de Darboux (\cref{exo:b1:derivative:10}), solo pueden
producir la última variedad: una función con una discontinuidad de
salto no es nunca la derivada de nada.
\end{example}

\begin{proposition}\label{prop:b1:continuity:algebra}
Las sumas, los productos, los cocientes (donde estén definidos) y las
composiciones de funciones \omterm{def:b1:continuity:continuous}{continuas} son \omterm{def:b1:continuity:continuous}{continuos}. Los \omterm{def:b1:poly:def}{polinomios},
las \omterm{def:b1:fractions:field}{fracciones racionales} (fuera de sus \omterm{def:b1:fractions:field}{polos}),
$\abs{\,\cdot\,}$, $\exp$, $\ln$, las funciones trigonométricas, las
\omterm{def:b1:functions:hyperbolic}{funciones hiperbólicas} y sus inversas (\cref{ch:b1:functions}) son
\omterm{def:b1:continuity:continuous}{continuas} en sus dominios.
\end{proposition}

\begin{proof}
Operaciones: \cref{cor:b1:continuity:operations}. Las constantes y la
identidad son \omterm{def:b1:continuity:continuous}{continuas} directamente por la definición
($\delta = \varepsilon$ sirve para la identidad y cualquier $\delta$
para las constantes); siendo \omterm{def:b1:continuity:continuous}{continuos} los productos de funciones
\omterm{def:b1:continuity:continuous}{continuas}, cada monomio $a_k x^k$ se sigue por inducción sobre $k$,
y las sumas rematan los \omterm{def:b1:poly:def}{polinomios}; una \omterm{def:b1:fractions:field}{fracción racional} es un
cociente de dos \omterm{def:b1:poly:def}{polinomios}, \omterm{def:b1:continuity:continuous}{continua} allí donde el denominador no se
anula. $\abs{\,\cdot\,}$: por la desigualdad triangular inversa,
$\bigl|\abs{f(x)} - \abs{f(x_0)}\bigr| \leq \abs{f(x) - f(x_0)}$,
luego sirve el mismo $\delta$. Para las funciones clásicas admitimos
aquí la \omterm{def:b1:continuity:continuous}{continuidad}; la derivabilidad (demostrada en
el~\cref{ch:b1:derivative}) es más fuerte.
\end{proof}

\begin{example}[Máximo y mínimo de funciones continuas]\label{ex:b1:continuity:maxmin}
Si $f, g$ son \omterm{def:b1:continuity:continuous}{continuas}, también lo son $\max(f, g)$ y
$\min(f, g)$: no hace falta distinguir casos, gracias a las identidades
\[
\max(f, g) = \frac{f + g + \abs{f - g}}{2},
\qquad
\min(f, g) = \frac{f + g - \abs{f - g}}{2},
\]
y a la \omterm{def:b1:continuity:continuous}{continuidad} de las sumas y de $\abs{\,\cdot\,}$
(\cref{prop:b1:continuity:algebra}). En particular,
$f^+ = \max(f, 0)$ y $f^- = \max(-f, 0)$ son \omterm{def:b1:continuity:continuous}{continuas} con
$f = f^+ - f^-$: la separación por signos que se usa en las series
(\cref{ch:b1:series}) y, a gran escala, en la teoría de la
integración del volumen del tercer año no cuesta nada en regularidad.
\end{example}

\begin{theorem}[Teorema del valor intermedio]\label{thm:b1:continuity:ivt}
Sea $f$ \omterm{def:b1:continuity:continuous}{continua} en $\intcc{a}{b}$ con $f(a) \leq 0 \leq f(b)$.
Entonces $f(c) = 0$ para algún $c \in \intcc{a}{b}$. En
consecuencia, una función \omterm{def:b1:continuity:continuous}{continua} en un \omterm{prop:b1:reals:intervals}{intervalo} toma todos los
valores comprendidos entre dos cualesquiera de sus valores: $f(I)$ es
un \omterm{prop:b1:reals:intervals}{intervalo}.\index{teorema del valor intermedio}
\end{theorem}

\begin{proof}
Por dicotomía. Póngase $a_0 = a$, $b_0 = b$. Dado $\intcc{a_n}{b_n}$
con $f(a_n) \leq 0 \leq f(b_n)$, sea $m$ el punto medio: si
$f(m) \leq 0$ consérvese $\intcc{m}{b_n}$, y si no,
$\intcc{a_n}{m}$; las condiciones de signo se mantienen. Las
sucesiones $(a_n), (b_n)$ son adyacentes
($b_n - a_n = \frac{b-a}{2^n}$), con límite común $c$
(\cref{thm:b1:seq:adjacent}). Por \omterm{def:b1:continuity:continuous}{continuidad} y por
el~\cref{thm:b1:seq:order}: $f(c) = \lim f(a_n) \leq 0$ y
$f(c) = \lim f(b_n) \geq 0$, luego $f(c) = 0$.

Para la consecuencia: dados valores $f(u) < v < f(w)$, aplíquese lo
anterior a $x \mapsto f(x) - v$ en el segmento de extremos $u$ y $w$;
así $f(I)$ es convexo, es decir, un \omterm{prop:b1:reals:intervals}{intervalo}
(\cref{prop:b1:reals:intervals}).
\end{proof}

\begin{omfigure}
\begin{tikzpicture}
\begin{axis}[omaxis, width=8.4cm, height=5.6cm,
    xmin=0, xmax=4.4, ymin=-1.1, ymax=1.4,
    xtick={0.6,3.8}, xticklabels={$a$,$b$}, ytick=\empty]
  \addplot[omDef, very thick, domain=0.6:3.8, samples=160]
    {0.5*(x-2) + 0.4*sin(deg(2.5*x))};
  \draw[dashed, gray] (axis cs:0.6,0) -- (axis cs:0.6,-0.301);
  \draw[dashed, gray] (axis cs:3.8,0) -- (axis cs:3.8,0.870);
  \node[below, font=\small, omDef] at (axis cs:0.6,-0.32) {$f(a)$};
  \node[above, font=\small, omDef] at (axis cs:3.8,0.89) {$f(b)$};
\end{axis}
\end{tikzpicture}

{\small Una función \omterm{def:b1:continuity:continuous}{continua} con $f(a) < 0 < f(b)$ tiene que cruzar
el eje: la demostración por dicotomía del~\cref{thm:b1:continuity:ivt}
atrapa un punto de cruce entre \omterm{thm:b1:seq:adjacent}{sucesiones adyacentes}.}
\end{omfigure}

\begin{example}[Una ecuación, el protocolo completo]\label{ex:b1:continuity:protocol}
Resuélvase $\eu^x = 3 - x$ en $\R$: existencia, unicidad y
localización. Póngase $g(x) = \eu^x + x - 3$, \omterm{def:b1:continuity:continuous}{continua}. Localización
y existencia: $g(0) = -2 < 0$ y $g(1) = \eu - 2 > 0$, de modo que el
teorema del valor intermedio planta una solución en $\intoo{0}{1}$.
Unicidad: $g$ es suma de $\eu^x$, estrictamente creciente, y de
$x - 3$, luego es estrictamente creciente en $\R$; y una función
estrictamente monótona toma cada valor a lo sumo una vez, así que la
solución es única en todo $\R$ (no solo en el \omterm{prop:b1:reals:intervals}{intervalo} sondeado). El
protocolo ---reescribir como $g = 0$, cambio de signo para la
existencia, monotonía para la unicidad--- resuelve en tres líneas casi
todas las preguntas de tipo «cuántas soluciones», y las tablas de
variación del~\cref{ch:b1:derivative} lo extienden a $g$ no monótonas
cortando $\R$ en ramas monótonas.
\end{example}

\begin{example}[La dicotomía como algoritmo]\label{ex:b1:continuity:dichotomy}
La demostración del~\cref{thm:b1:continuity:ivt} calcula. Tómese
$f(x) = x^3 + x - 1$: $f(0) = -1 < 0 < 1 = f(1)$, luego hay una raíz
en $\intoo{0}{1}$. Bisecando:
\[
f(0.5) = -0.375 < 0, \qquad
f(0.75) = 0.171875 > 0, \qquad
f(0.625) = -0.130859375 < 0 ,
\]
así que la raíz queda atrapada sucesivamente en $\intoo{0.5}{1}$,
después en $\intoo{0.5}{0.75}$ y después en $\intoo{0.625}{0.75}$
(valor verdadero: $c \approx 0.6823$). Tras $n$ pasos, el error es a
lo sumo $\frac{b - a}{2^n}$: diez pasos dan tres decimales, y veinte,
seis. La idea de cierre: el teorema del valor intermedio no es solo un
\omterm{def:b1:logic:statement}{enunciado} de existencia --- su demostración por dicotomía es un
algoritmo garantizado, aunque lento, de búsqueda de raíces, frente al
cual conviene medir el rápido pero local método de Newton
del~\cref{ch:b1:derivative}.
\end{example}

\begin{theorem}[Teorema de los valores extremos]\label{thm:b1:continuity:evt}
Una función \omterm{def:b1:continuity:continuous}{continua} en un segmento $\intcc{a}{b}$ está acotada y
alcanza sus cotas: existen $c, d \in \intcc{a}{b}$ con
\[
f(c) = \inf_{\intcc{a}{b}} f,
\qquad
f(d) = \sup_{\intcc{a}{b}} f .
\]
Combinado con el~\cref{thm:b1:continuity:ivt}: \emph{la imagen
\omterm{def:b1:continuity:continuous}{continua} de un segmento es un segmento} $\intcc{f(c)}{f(d)}$.
\index{teorema de los valores extremos}
\end{theorem}

\begin{proof}
\emph{Acotada superiormente:} si no, tómense $u_n$ con
$f(u_n) \geq n$. Por la compacidad del segmento
(\cref{thm:b1:topology:compact}), una \omterm{def:b1:seq:subsequence}{subsucesión}
$u_{\varphi(n)} \to x \in \intcc{a}{b}$; la \omterm{def:b1:continuity:continuous}{continuidad} da
$f(u_{\varphi(n)}) \to f(x)$, pero
$f(u_{\varphi(n)}) \geq \varphi(n) \to +\infty$: contradicción.

\emph{El \omterm{def:b1:reals:bounds}{supremo} se alcanza:} sea $M = \sup f$ y elíjanse $v_n$ con
$f(v_n) > M - \frac{1}{n+1}$ (\cref{prop:b1:reals:epsilon}).
Extráigase $v_{\varphi(n)} \to d \in \intcc{a}{b}$: entonces
$f(d) = \lim f(v_{\varphi(n)}) = M$ por emparedado. El \omterm{def:b1:reals:bounds}{ínfimo} se
trata con $-f$.
\end{proof}

\begin{example}[Mínimo positivo en un segmento]\label{ex:b1:continuity:posmin}
Sea $f$ \omterm{def:b1:continuity:continuous}{continua} en $\intcc{0}{1}$ con $f(x) > 0$ para todo $x$.
Entonces $\inf f > 0$: por el teorema de los valores extremos, el
\omterm{def:b1:reals:bounds}{ínfimo} es un \emph{valor} $f(c)$, y $f(c) > 0$ por hipótesis. Así
pues, una función \omterm{def:b1:continuity:continuous}{continua} y positiva en un segmento se mantiene
lejos de $0$ --- un argumento de dos líneas que se usará una docena
de veces en los capítulos siguientes (denominadores bajo control,
encuadres por funciones escalonadas, cotas de error). En un
\omterm{prop:b1:reals:intervals}{intervalo} no compacto esto falla espectacularmente: $f(x) = x$ en
$\intoc{0}{1}$ es \omterm{def:b1:continuity:continuous}{continua} y positiva con $\inf f = 0$, no alcanzado.
La idea de cierre: «positiva» asciende a «uniformemente positiva»
exactamente cuando el dominio es compacto; toda hipótesis del teorema
de los valores extremos sostiene el edificio.
\end{example}

\begin{remark}[Errores frecuentes en torno a los tres teoremas]
(i) \emph{Imágenes \omterm{def:b1:continuity:continuous}{continuas}}: solo los \emph{segmentos} son robustos.
La imagen \omterm{def:b1:continuity:continuous}{continua} de un \omterm{prop:b1:reals:intervals}{intervalo} \omterm{def:b1:topology:open}{abierto} no tiene por qué ser
\omterm{def:b1:topology:open}{abierta} ($x \mapsto x^2$ aplica $\intoo{-1}{1}$ sobre
$\intco{0}{1}$), ni la imagen de un \omterm{def:b1:topology:closed}{cerrado} tiene por qué ser
\omterm{def:b1:topology:closed}{cerrada} ($\arctan$ aplica el \omterm{def:b1:topology:closed}{cerrado} $\R$ sobre el \omterm{def:b1:topology:open}{abierto}
$\intoo{-\frac\pi2}{\frac\pi2}$); pero la imagen de un segmento es un
segmento (\cref{thm:b1:continuity:evt}). (ii) \emph{El teorema del
valor intermedio necesita un \omterm{prop:b1:reals:intervals}{intervalo}}: la función
$x \mapsto \frac1x$, \omterm{def:b1:continuity:continuous}{continua} en
$\intco{-1}{0} \cup \intoc{0}{1}$, toma los valores $-1$ y $1$ y, sin
embargo, no se anula nunca --- su dominio son dos piezas disjuntas y
el valor $0$ cae en el hueco; nómbrese siempre el \omterm{prop:b1:reals:intervals}{intervalo} sobre el
que se aplica el teorema. (iii) \emph{La \omterm{def:b1:continuity:uniform}{continuidad uniforme} es una
propiedad de la pareja (función, \omterm{def:b1:logic:sets}{conjunto})}: $x^2$ es
\omterm{def:b1:continuity:uniform}{uniformemente continua} en todo segmento pero no en $\R$
(\cref{ex:b1:continuity:notuniform}) --- las palabras
«\omterm{def:b1:continuity:uniform}{uniformemente continua}» sin un dominio no significan nada. (iv)
\emph{La \omterm{def:b1:continuity:continuous}{continuidad} de la inversa no es formal}: se cumple en
\omterm{prop:b1:reals:intervals}{intervalos} gracias a la monotonía
(\cref{thm:b1:continuity:bijection}), pero una \omterm{def:b1:logic:inj}{biyección} \omterm{def:b1:continuity:continuous}{continua}
entre uniones de \omterm{prop:b1:reals:intervals}{intervalos} puede tener inversa discontinua --- la
observación que sigue a ese teorema está ahí porque se cita sin la
hipótesis de \omterm{prop:b1:reals:intervals}{intervalo}.
\end{remark}

\section{Funciones monótonas y funciones inversas}

\begin{theorem}[Teorema de la biyección monótona]\label{thm:b1:continuity:bijection}
Sea $f$ \omterm{def:b1:continuity:continuous}{continua} y estrictamente monótona en un \omterm{prop:b1:reals:intervals}{intervalo} $I$. Entonces:
\begin{enumerate}
  \item $f$ es una \omterm{def:b1:logic:inj}{biyección} de $I$ sobre el \omterm{prop:b1:reals:intervals}{intervalo} $J = f(I)$;
  \item la inversa $f^{-1} \colon J \to I$ es estrictamente monótona
        (en el mismo sentido) y \emph{\omterm{def:b1:continuity:continuous}{continua}}.
\end{enumerate}
\end{theorem}

\begin{proof}
Digamos que $f$ es estrictamente creciente. (1) La \omterm{def:b1:logic:inj}{inyectividad} es
inmediata por la monotonía estricta; la \omterm{def:b1:logic:inj}{sobreyectividad} sobre $f(I)$
es trivial, y $f(I)$ es un \omterm{prop:b1:reals:intervals}{intervalo} por el~\cref{thm:b1:continuity:ivt}.

(2) $f^{-1}$ es estrictamente creciente: si $y < y'$ en $J$ pero
$f^{-1}(y) \geq f^{-1}(y')$, aplicando la creciente $f$ se obtiene
$y \geq y'$, absurdo. \omterm{def:b1:continuity:continuous}{Continuidad} de $f^{-1}$ en
$y_0 = f(x_0) \in J$: sea $\varepsilon > 0$. Supóngase primero $x_0$
\omterm{def:b1:topology:closure}{interior} a $I$, y encójase $\varepsilon$ de modo que
$x_0 \pm \varepsilon \in I$: sus imágenes cumplen
$f(x_0 - \varepsilon) < y_0 < f(x_0 + \varepsilon)$. Tómese
$\delta = \min\bigl(y_0 - f(x_0 - \varepsilon),\, f(x_0 +
\varepsilon) - y_0\bigr) > 0$: para $\abs{y - y_0} \leq \delta$, la
monotonía de $f^{-1}$ aprieta $f^{-1}(y)$ entre $x_0 - \varepsilon$ y
$x_0 + \varepsilon$. Si $x_0$ es, digamos, el extremo izquierdo de
$I$, solo se dispone de $x_0 + \varepsilon$: entonces
$y_0 = \min J$ (monotonía), todo $y \in J$ con
$y - y_0 \leq f(x_0 + \varepsilon) - y_0$ cumple
$x_0 \leq f^{-1}(y) \leq x_0 + \varepsilon$, y la estimación
unilateral es exactamente la \omterm{def:b1:continuity:continuous}{continuidad} en un extremo; el extremo
derecho es simétrico. (Nótese que la \omterm{def:b1:continuity:continuous}{continuidad} de $f^{-1}$
\emph{no} se ha deducido de la de $f$ por simetría --- lo que hace el
trabajo es la monotonía sobre un \omterm{prop:b1:reals:intervals}{intervalo}.)
\end{proof}

\begin{remark}
Este teorema es lo que la~\cref{def:b1:functions:arc} y
la~\cref{prop:b1:functions:invhyp} usaron en silencio: $\arcsin$,
$\arccos$, $\arctan$, $\operatorname{arsinh}$, \dots\ son
\omterm{def:b1:continuity:continuous}{continuas}. Un complemento (\cref{exo:b1:continuity:10}): una
función \omterm{def:b1:continuity:continuous}{continua} e \emph{\omterm{def:b1:logic:inj}{inyectiva}} en un \omterm{prop:b1:reals:intervals}{intervalo} es
automáticamente estrictamente monótona, de modo que la hipótesis de
monotonía no cuesta nada.
\end{remark}

\begin{example}[Raíces de todos los órdenes]\label{ex:b1:continuity:nthroot}
Para $n \in \N^*$, la función $f(x) = x^n$ es \omterm{def:b1:continuity:continuous}{continua} y
estrictamente creciente en $\intco{0}{+\infty}$, con $f(0) = 0$ y
$f(x) \to +\infty$: su imagen es $\intco{0}{+\infty}$
(\cref{thm:b1:continuity:ivt} para la estructura de \omterm{prop:b1:reals:intervals}{intervalo}). El
teorema de la \omterm{def:b1:logic:inj}{biyección} monótona entrega entonces, de un solo golpe,
una inversa \emph{\omterm{def:b1:continuity:continuous}{continua}} y estrictamente creciente
\[
x \mapsto x^{1/n} \colon \intco{0}{+\infty} \to
\intco{0}{+\infty} :
\]
la existencia, la unicidad y la \omterm{def:b1:continuity:continuous}{continuidad} de las raíces $n$-ésimas,
sin ningún cálculo. Compárese con el~\cref{exo:b1:reals:12}, que
construía $\sqrt y$ a mano a partir del \omterm{def:b1:reals:bounds}{supremo}: un capítulo de teoría
ha comprimido esa página de trabajo en dos líneas, y esas mismas dos
líneas legitimaron antes $\arcsin$, $\arctan$ y
$\operatorname{arsinh}$. La idea de cierre: un buen teorema es trabajo
almacenado.
\end{example}

\section{Continuidad uniforme}

\begin{definition}\label{def:b1:continuity:uniform}
$f \colon I \to \R$ es \emph{uniformemente continua}\index{continuidad
uniforme} cuando
\[
\forall \varepsilon > 0,\ \exists \delta > 0,\ \forall x, y \in I,
\qquad \abs{x - y} \leq \delta \implies \abs{f(x) - f(y)} \leq
\varepsilon .
\]
La clave: $\delta$ depende solo de $\varepsilon$, no del lugar de
$I$. La \omterm{def:b1:continuity:continuous}{continuidad} uniforme implica la \omterm{def:b1:continuity:continuous}{continuidad}; y una función
\emph{lipschitziana} ($\abs{f(x) - f(y)} \leq k \abs{x - y}$) es
uniformemente \omterm{def:b1:continuity:continuous}{continua} ($\delta = \varepsilon/k$).
\end{definition}

\begin{example}\label{ex:b1:continuity:notuniform}
$x \mapsto x^2$ es \omterm{def:b1:continuity:continuous}{continua} en $\R$ pero \emph{no}
\omterm{def:b1:continuity:uniform}{uniformemente continua}:
$\abs{(n + \frac 1n)^2 - n^2} = 2 + \frac{1}{n^2} \geq 2$ aunque los
argumentos están a distancia $\frac 1n \to 0$. En cualquier \omterm{prop:b1:reals:intervals}{intervalo}
acotado es lipschitziana y, por tanto, \omterm{def:b1:continuity:uniform}{uniformemente continua} ---
en consonancia con el teorema de Heine de más abajo.
\end{example}

\begin{example}[Módulos uniformes, explícitamente]\label{ex:b1:continuity:moduli}
En un segmento, Heine garantiza un $\delta$ uniforme; y a menudo
también se puede \emph{calcular}. Para $f(x) = x^2$ en $\intcc{0}{10}$:
\[
\abs{x^2 - y^2} = \abs{x + y}\,\abs{x - y} \leq 20\,\abs{x - y},
\]
así que $\delta = \frac{\varepsilon}{20}$ sirve uniformemente (un
\omterm{def:b1:complex:field}{módulo} lipschitziano, lineal en $\varepsilon$). Para $\sqrt x$ en
$\intcc{0}{1}$: $\abs{\sqrt x - \sqrt y} \leq \sqrt{\abs{x - y}}$
(\cref{exo:b1:continuity:9}), luego sirve $\delta = \varepsilon^2$ ---
uniforme pero \emph{no} lineal: cerca de $0$ la raíz cuadrada es
empinada, y el precio aparece en el exponente de $\varepsilon$, no en
un fallo de la uniformidad. La idea de cierre: la
\omterm{def:b1:continuity:uniform}{continuidad uniforme} es un espectro, no un sí/no --- la función
$\delta(\varepsilon)$, llamada \omterm{def:b1:complex:field}{módulo}, mide lo cara que es la
uniformidad, y ser lipschitziana es sencillamente su mejor nota.
\end{example}

\begin{theorem}[Heine]\label{thm:b1:continuity:heine}
Una función \omterm{def:b1:continuity:continuous}{continua} en un \emph{segmento} $\intcc{a}{b}$ es
\omterm{def:b1:continuity:uniform}{uniformemente continua}.\index{teorema de Heine}
\end{theorem}

\begin{proof}
Por reducción al absurdo: supóngase que algún $\varepsilon_0 > 0$
derrota a todo $\delta$. Con $\delta = \frac{1}{n+1}$, tómense
$x_n, y_n \in \intcc{a}{b}$ con $\abs{x_n - y_n} \leq \frac{1}{n+1}$
y $\abs{f(x_n) - f(y_n)} > \varepsilon_0$. La compacidad
(\cref{thm:b1:topology:compact}) extrae
$x_{\varphi(n)} \to c \in \intcc{a}{b}$; y entonces
$y_{\varphi(n)} \to c$ también (emparedado sobre $\abs{x - y}$). La
\omterm{def:b1:continuity:continuous}{continuidad} en $c$ da $f(x_{\varphi(n)}) \to f(c)$ y
$f(y_{\varphi(n)}) \to f(c)$, luego
$\abs{f(x_{\varphi(n)}) - f(y_{\varphi(n)})} \to 0 < \varepsilon_0$:
contradicción.
\end{proof}

\begin{example}[Acotada, continua y, aun así, no uniformemente]\label{ex:b1:continuity:sinx2}
La función $f(x) = \sin(x^2)$ es \omterm{def:b1:continuity:continuous}{continua} y acotada en $\R$, pero no
\omterm{def:b1:continuity:uniform}{uniformemente continua}. Tómense
\[
x_n = \sqrt{2\pi n}, \qquad y_n = \sqrt{2\pi n + \tfrac\pi2}:
\qquad
y_n - x_n = \frac{\pi/2}{\sqrt{2\pi n + \frac\pi2} +
\sqrt{2\pi n}} \longrightarrow 0 ,
\]
y, sin embargo,
$f(y_n) - f(x_n) = \sin\bigl(2\pi n + \frac\pi2\bigr) - \sin(2\pi n)
= 1 - 0 = 1$ para todo $n$: ningún $\delta$ único puede servir para
$\varepsilon = \frac12$ en todas partes. Geométricamente, las
oscilaciones de $\sin(x^2)$ se \emph{aceleran}: la gráfica completa
una onda entera en ventanas cada vez más cortas, de modo que la
escala horizontal que exige un $\varepsilon$ dado se encoge hasta
cero conforme $x$ crece. La idea de cierre: estar acotada no compra
uniformidad (este ejemplo), y no estarlo no la impide ($\sqrt x$,
\cref{exo:b1:continuity:9}); lo que decide es el \emph{módulo de
oscilación}, y el teorema de Heine dice que los dominios compactos lo
disciplinan automáticamente.
\end{example}

\begin{example}[El experimento de la calculadora, explicado]\label{ex:b1:continuity:cositeration}
Tecléese un número cualquiera en una calculadora y púlsese $\cos$
repetidamente: la pantalla se ancla en $0.7390851\dots$ ¿Por qué?
Tras una pulsación el valor está en $\intcc{-1}{1}$; tras dos, en
$\intcc{\cos 1}{1} \subseteq \intcc{0.54}{1}$, un \omterm{prop:b1:reals:intervals}{intervalo} estable
para $\cos$. En él, $\abs{\cos a - \cos b} \leq \sin(1)\,\abs{a - b}$
con $\sin 1 = 0.841\dots < 1$ (la cota de producto a suma
del~\cref{pb:b1:seq:1}, o la desigualdad del valor medio
del~\cref{ch:b1:derivative}): la iteración contrae, luego por el paso
de control del error del~\cref{met:b1:seq:recurrent},
\[
\abs{u_n - c} \leq (0.842)^{\,n-2}\,\abs{u_2 - c}
\longrightarrow 0 ,
\]
donde $c$ es el único punto fijo $\cos c = c$
(\cref{exo:b1:continuity:6}). Unas $40$ \omterm{pb:b1:diffeq:1}{pulsaciones} compran tres
decimales ($0.842^{40} \approx 10^{-3}$) --- un ritmo geométrico, más
suave que el $2^{-n}$ por paso de la dicotomía, pero cada pulsación
cuesta una tecla mientras que cada paso de dicotomía cuesta una
evaluación completa de signo. La idea de cierre: la imagen del punto
fijo del~\cref{ch:b1:seq} y los teoremas de existencia de este
capítulo son las dos mitades de una misma historia --- el teorema del
valor intermedio encuentra $c$, y la contracción llega hasta él.
\end{example}

\begin{remark}[Dónde funcionan estos teoremas a continuación]
Cada uno de los tres pilares de este capítulo alimenta un capítulo
posterior. El teorema del valor intermedio nutre todo argumento de
existencia de soluciones y el teorema de la \omterm{def:b1:logic:inj}{biyección} monótona; el de
los valores extremos convierte los problemas de optimización en
teoremas (el de Rolle y el del valor medio del~\cref{ch:b1:derivative}
arrancan justo ahí); y el de Heine es la razón de que las funciones
\omterm{def:b1:continuity:continuous}{continuas} en segmentos se puedan integrar en
el~\cref{ch:b1:integration} --- el $\delta$ uniforme es lo que hace
converger las sumas de Riemann. En el volumen del segundo año, el
mismo trío reaparece en los espacios vectoriales normados, con la
compacidad haciendo el trabajo que aquí hacen los segmentos.
\end{remark}

\begin{remark}[Perspectivas dentro de este volumen]
La \omterm{def:b1:continuity:continuous}{continuidad} está a punto de ser superada, pero nunca jubilada.
\Cref{ch:b1:derivative} la refuerza hasta la derivabilidad y le
devuelve el favor (derivable implica \omterm{def:b1:continuity:continuous}{continua});
el~\cref{ch:b1:integration} descansa dos veces sobre ella, por medio
de Heine para la construcción y por medio del teorema fundamental,
cuyo objeto central $x \mapsto \int_a^x f$ eleva una $f$ meramente
\omterm{def:b1:continuity:continuous}{continua} a una primitiva de clase $C^1$. En
el~\cref{ch:b1:multivar}, la \omterm{def:b1:continuity:continuous}{continuidad} en dos variables esconde una
trampa que merece anticiparse: la función $\frac{xy}{x^2 + y^2}$
(extendida por $0$) es \omterm{def:b1:continuity:continuous}{continua} en $x$ para cada $y$ fijo y en $y$
para cada $x$ fijo y, aun así, no es \omterm{def:b1:continuity:continuous}{continua} en el origen --- a lo
largo de la diagonal $x = y$ vale constantemente $\frac12$. La
\omterm{def:b1:continuity:continuous}{continuidad} separada es estrictamente más débil que la
\omterm{def:b1:continuity:continuous}{continuidad}: la caracterización secuencial sobrevive al salto a
$\R^2$, pero hay que permitir que las sucesiones se acerquen desde
\emph{todas} las direcciones, no solo a lo largo de los ejes.
\end{remark}

\section{Ejercicios}

\begin{exercise}[$\star$]\label{exo:b1:continuity:1}
Usando la caracterización secuencial, demuéstrese que
$x \mapsto \sin\frac 1x$ no tiene límite en $0^+$ \emph{(exhíbanse dos
sucesiones)}. ¿Lo tiene $x \mapsto x \sin\frac 1x$?
\end{exercise}

\begin{exercise}[$\star$]\label{exo:b1:continuity:2}
Estúdiese la \omterm{def:b1:continuity:continuous}{continuidad} en $\R$ de $f(x) = \lfloor x \rfloor$, de
$g(x) = x - \lfloor x \rfloor$ y de
$h(x) = \lfloor x \rfloor + (x - \lfloor x\rfloor)^2$.
\end{exercise}

\begin{exercise}[$\star$]\label{exo:b1:continuity:3}
Demuéstrese que la ecuación $x^5 - 3x + 1 = 0$ tiene al menos tres
soluciones reales \emph{(evalúese en puntos bien elegidos y aplíquese
el~\cref{thm:b1:continuity:ivt} en tres segmentos disjuntos)}.
\end{exercise}

\begin{exercise}[$\star$]\label{exo:b1:continuity:4}
Demuéstrese que todo \omterm{def:b1:poly:def}{polinomio} de grado impar tiene una raíz real.
\end{exercise}

\begin{exercise}[$\star\star$]\label{exo:b1:continuity:5}
(Punto fijo) Sea $f \colon \intcc{0}{1} \to \intcc{0}{1}$
\omterm{def:b1:continuity:continuous}{continua}. Demuéstrese que $f$ tiene un punto fijo: $f(c) = c$ para
algún $c$. Ilústrese que no se puede prescindir ni de la
\omterm{def:b1:continuity:continuous}{continuidad} ni del segmento.
\end{exercise}

\begin{exercise}[$\star\star$]\label{exo:b1:continuity:6}
Demuéstrese que la ecuación $\cos x = x$ tiene exactamente una
solución real y que está en $\intoo{0}{1}$.
\end{exercise}

\begin{exercise}[$\star\star$]\label{exo:b1:continuity:7}
Sea $f \colon \R \to \R$ \omterm{def:b1:continuity:continuous}{continua} con $f(x) \to +\infty$ cuando
$x \to \pm\infty$. Demuéstrese que $f$ alcanza un mínimo global en
$\R$. \emph{(Redúzcase a un segmento que contenga un \omterm{def:b1:logic:sets}{conjunto} de subnivel.)}
\end{exercise}

\begin{exercise}[$\star\star$]\label{exo:b1:continuity:8}
Sea $f \colon \R \to \R$ \omterm{def:b1:continuity:continuous}{continua} y periódica (de período
$T > 0$). Demuéstrese que $f$ está acotada y alcanza sus cotas, y que
existe $c$ con $f(c + \frac T2) = f(c)$. \emph{(Para lo segundo,
estúdiese $g(x) = f(x + \frac T2) - f(x)$ sobre un período.)}
\end{exercise}

\begin{exercise}[$\star\star$]\label{exo:b1:continuity:9}
Demuéstrese que $x \mapsto \sqrt x$ es \omterm{def:b1:continuity:uniform}{uniformemente continua} en
$\intco{0}{+\infty}$, aunque no sea lipschitziana cerca de $0$.
\emph{(Demuéstrese y úsese $\abs{\sqrt x - \sqrt y} \leq \sqrt{\abs{x -
y}}$.)}
\end{exercise}

\begin{exercise}[$\star\star\star$]\label{exo:b1:continuity:10}
Sea $f$ \omterm{def:b1:continuity:continuous}{continua} e \omterm{def:b1:logic:inj}{inyectiva} en un \omterm{prop:b1:reals:intervals}{intervalo} $I$. Demuéstrese que
$f$ es estrictamente monótona. \emph{Indicación: si no, hay
$a < b < c$ con, digamos, $f(b) > f(a)$ y $f(b) > f(c)$; aplíquese el
teorema del valor intermedio a un valor entre $\max(f(a), f(c))$ y
$f(b)$ a ambos lados de $b$.}
\end{exercise}

\begin{exercise}[$\star\star\star$]\label{exo:b1:continuity:11}
(Ecuación funcional de Cauchy, caso \omterm{def:b1:continuity:continuous}{continuo}) Sea
$f \colon \R \to \R$ \omterm{def:b1:continuity:continuous}{continua} con $f(x + y) = f(x) + f(y)$ para
todos $x, y$. Demuéstrese que $f(x) = f(1)\,x$ para todo $x$: primero
en $\N$, $\Z$ y $\Q$ (solo con la aditividad), y después en $\R$ por
\omterm{def:b1:continuity:continuous}{continuidad} y \omterm{def:b1:topology:dense}{densidad} (\cref{thm:b1:reals:density}).
\end{exercise}

\begin{exercise}[$\star\star\star$]\label{exo:b1:continuity:12}
Sea $f \colon \intco{0}{+\infty} \to \R$ \omterm{def:b1:continuity:continuous}{continua} con
$f(x) \to \ell \in \R$ cuando $x \to +\infty$. Demuéstrese que $f$ es
\emph{uniformemente} \omterm{def:b1:continuity:continuous}{continua} en $\intco{0}{+\infty}$. \emph{(Córtese
en un $A$ grande: Heine en $\intcc{0}{A+1}$ y el límite más allá de
$A$; háganse solapar los dos regímenes.)}
\end{exercise}

\section{Problema: la ecuación funcional de Cauchy y sus hermanas}

\begin{problem}[{Problema del fin de semana --- $f(x + y) = f(x) + f(y)$:
la regularidad fuerza la linealidad, y el retrato de los monstruos}]
\label{pb:b1:continuity:1}
¿Qué funciones cumplen $f(x + y) = f(x) + f(y)$ para todos los reales
$x, y$? Cauchy planteó la pregunta en 1821; la respuesta es un
paradigma. \Cref{exo:b1:continuity:11} muestra que una función así,
llamada \emph{aditiva}, es lineal en $\Q$ y que la \omterm{def:b1:continuity:continuous}{continuidad} plena
fuerza $f(x) = cx$. Este problema afina la hipótesis de forma
drástica ---la \omterm{def:b1:continuity:continuous}{continuidad} en un \emph{solo} punto, o la monotonía,
o la mera acotación en un pequeño \omterm{prop:b1:reals:intervals}{intervalo}, bastan cada una por sí
sola---, después pinta el retrato de una hipotética solución no
lineal (su gráfica llena el plano), resuelve las ecuaciones hermanas
que caracterizan $\eu^{cx}$, $c\ln x$, $x^c$ y $cx^2$, y cierra con la
ecuación de Jensen y el teorema \emph{convexa en el punto medio $+$
\omterm{def:b1:continuity:continuous}{continua} $\implies$ convexa}.\index{ecuación funcional de
Cauchy}\index{ecuación funcional} En todo el problema, \emph{aditiva}
significa: $f(x + y) = f(x) + f(y)$ para todos $x, y \in \R$.

\medskip
\noindent\textbf{Parte I --- $\Q$-linealidad, y un punto de
\omterm{def:b1:continuity:continuous}{continuidad}.}
\begin{enumerate}
  \item Sea $f$ aditiva. Por el~\cref{exo:b1:continuity:11},
        $f(r) = f(1)\,r$ para $r$ racional. Demuéstrese el
        \omterm{def:b1:logic:statement}{enunciado} más fino que se usará más abajo: para \emph{todo}
        $x \in \R$ y todo $r \in \Q$, $f(rx) = r\,f(x)$ ($f$ es
        $\Q$-\emph{lineal}).
  \item Supóngase que la $f$ aditiva es \omterm{def:b1:continuity:continuous}{continua} en un solo punto
        $x_0$. Véase que $f$ es \omterm{def:b1:continuity:continuous}{continua} en todas partes
        \emph{(calcúlese $f(x + h) - f(x)$ en función de
        $f(x_0 + h) - f(x_0)$)}, luego $f(x) = f(1)\,x$.
  \item Véase que una función aditiva queda determinada por su
        restricción a cualquier \omterm{def:b1:structures:subgroup}{subgrupo} \omterm{def:b1:topology:dense}{denso}: si dos funciones
        aditivas coinciden en $\Z + \sqrt2\,\Z$ y las dos son
        \omterm{def:b1:continuity:continuous}{continuas}, son iguales --- mientras que, sin \omterm{def:b1:continuity:continuous}{continuidad},
        prescribir $f(1) = 0$ y $f(\sqrt 2) = 1$ es compatible con la
        $\Q$-linealidad sobre el \omterm{def:b1:structures:subgroup}{subgrupo}. Calcúlese
        $f(m + n\sqrt2)$ para esta prescripción.
  \item Sea $f$ aditiva y acotada superiormente por $M$ en algún
        \omterm{prop:b1:reals:intervals}{intervalo} $\intcc{a}{b}$ con $a < b$. Véase que $f$ está
        acotada superiormente en $\intcc{0}{\ell}$, con
        $\ell = b - a$ \emph{(trasládese por $a$)}.
\end{enumerate}

\medskip
\noindent\textbf{Parte II --- La escalera de regularidad.}
\begin{enumerate}[resume]
  \item Continuando la pregunta 4: usando
        $f(\ell - t) + f(t) = f(\ell)$, véase que $f$ también está
        acotada \emph{inferiormente} en $\intcc{0}{\ell}$:
        $\abs f \leq C$ allí.
  \item Véase que $\abs{f(t)} \leq \frac{C}{n}$ para
        $t \in \intcc{0}{\ell/n}$, y dedúzcase que $f$ es \omterm{def:b1:continuity:continuous}{continua}
        en $0$ \emph{(la imparidad se ocupa del lado izquierdo)}, y
        por tanto en todas partes (pregunta 2): una función aditiva
        acotada en un \omterm{prop:b1:reals:intervals}{intervalo} es lineal.
  \item Dedúzcase el caso monótono: una función aditiva no
        decreciente en algún $\intcc{a}{b}$ ($a < b$) es $f(x) = cx$
        con $c \geq 0$.
  \item Móntese la \emph{escalera de regularidad}: para una $f$
        aditiva, son equivalentes --- (a) $f(x) = cx$; (b) $f$
        \omterm{def:b1:continuity:continuous}{continua}; (c) $f$ \omterm{def:b1:continuity:continuous}{continua} en un punto; (d) $f$ monótona
        en algún \omterm{prop:b1:reals:intervals}{intervalo} no degenerado; (e) $f$ acotada en algún
        \omterm{prop:b1:reals:intervals}{intervalo} no degenerado. Ordénense las implicaciones de modo
        que cada una sea trivial o ya esté demostrada.
  \item Compruébese que las preguntas 4--6 solo consumieron una cota
        \emph{superior}: una función aditiva acotada superiormente en
        un \omterm{prop:b1:reals:intervals}{intervalo} no degenerado ya es lineal. Dedúzcase el
        \omterm{def:b1:logic:statement}{enunciado} espejo para una cota inferior, y regístrese la
        forma más fuerte del peldaño (e) así obtenida.
\end{enumerate}

\medskip
\noindent\textbf{Parte III --- Retrato de un monstruo.} Supóngase
ahora $f$ aditiva pero \emph{no} lineal.
\begin{enumerate}[resume]
  \item Véase que hay reales no nulos $u, v$ con
        $\dfrac{f(u)}{u} \neq \dfrac{f(v)}{v}$, y que los vectores
        $(u, f(u))$ y $(v, f(v))$ generan el plano (su determinante
        $u f(v) - v f(u)$ es no nulo).
  \item Véase que la gráfica de $f$ contiene todos los puntos
        \[
        r\,(u, f(u)) + s\,(v, f(v)), \qquad r, s \in \Q ,
        \]
        y dedúzcase que la gráfica es \emph{\omterm{def:b1:topology:dense}{densa} en $\R^2$}: para
        todo punto $(x_0, y_0)$ del plano y todo $\varepsilon > 0$,
        algún $(x, f(x))$ dista de él menos de $\varepsilon$
        \emph{(resuélvase el sistema real $2\times2$ y aproxímense
        después los coeficientes reales por racionales)}.
  \item Dedúzcase de la pregunta 11 el retrato completo: una función
        aditiva no lineal no está acotada en ningún \omterm{prop:b1:reals:intervals}{intervalo} no
        degenerado, es discontinua en todo punto, no es monótona en
        ningún \omterm{prop:b1:reals:intervals}{intervalo} y su imagen de cualquier \omterm{prop:b1:reals:intervals}{intervalo} es \omterm{def:b1:topology:dense}{densa}
        en $\R$. Concíliese con la pregunta 8.
  \item Los monstruos existen --- sobre un \omterm{def:b1:structures:subgroup}{subgrupo} \omterm{def:b1:topology:dense}{denso}, y de forma
        constructiva: sobre $G = \Z + \sqrt2\,\Z$ defínase
        $f(m + n\sqrt2) = n$. Véase que $f$ está bien definida y es
        aditiva en $G$, y que $f$ no está acotada en
        $G \cap \intoo{0}{\varepsilon}$ para ningún $\varepsilon > 0$
        \emph{(para $N$ fijo, solo finitos $g = m + n\sqrt2 \in
        \intoo{0}{1}$ cumplen $\abs n \leq N$; y, en cambio,
        $G \cap \intoo{0}{\varepsilon}$ es infinito)}. Explíquese en
        un párrafo por qué extender una $f$ así a todo $\R$ exige una
        base de $\R$ como espacio vectorial sobre $\Q$ (una
        \emph{base de Hamel}), cuya existencia es un asunto del axioma
        de elección que queda fuera de este volumen.
\end{enumerate}

\medskip
\noindent\textbf{Parte IV --- Las ecuaciones hermanas.} Todas las
funciones de esta parte son \omterm{def:b1:continuity:continuous}{continuas}.
\begin{enumerate}[resume]
  \item Sea $f \colon \R \to \R$ \omterm{def:b1:continuity:continuous}{continua}, no idénticamente $0$, con
        $f(x + y) = f(x)f(y)$. Véase que
        $f(x) = f\bigl(\frac x2\bigr)^2 \geq 0$, después que $f > 0$
        en todas partes y después que $f(x) = \eu^{cx}$ para algún
        $c$: las exponenciales son exactamente los \omterm{def:b1:structures:morphism}{morfismos}
        \omterm{def:b1:continuity:continuous}{continuos} de $(\R, +)$ en $(\R^*, \times)$.
  \item Sea $f \colon \intoo{0}{+\infty} \to \R$ \omterm{def:b1:continuity:continuous}{continua} con
        $f(xy) = f(x) + f(y)$. Véase que $f(x) = c\ln x$
        \emph{(transpórtese por $\exp$)}.
  \item Sea $f \colon \intoo{0}{+\infty} \to \intoo{0}{+\infty}$
        \omterm{def:b1:continuity:continuous}{continua} con $f(xy) = f(x)f(y)$. Véase que $f(x) = x^c$.
  \item Hállense todas las $f \colon \R \to \R$ \omterm{def:b1:continuity:continuous}{continuas} con
        \[
        f(x + y) = f(x) + f(y) + f(x)f(y) .
        \]
        \emph{(Estúdiese $h = 1 + f$; trátese aparte el caso
        degenerado.)}
  \item (Paralelogramo) Hállense todas las $f \colon \R \to \R$
        \omterm{def:b1:continuity:continuous}{continuas} con $f(x + y) + f(x - y) = 2f(x) + 2f(y)$:
        véase que $f$ es par, que $f(0) = 0$, que
        $f(nx) = n^2 f(x)$ por inducción, y después que
        $f(x) = f(1)\,x^2$. (Esta ecuación es la huella dactilar de
        las formas cuadráticas --- la ley del paralelogramo que
        detecta, en el volumen del segundo año, qué normas provienen
        de un producto escalar.)
\end{enumerate}

\medskip
\noindent\textbf{Parte V --- Jensen y la convexidad en el punto medio.}
\begin{enumerate}[resume]
  \item (Ecuación de Jensen) Sea $f \colon \R \to \R$ \omterm{def:b1:continuity:continuous}{continua} con
        $f\bigl(\frac{x+y}{2}\bigr) = \frac{f(x) + f(y)}{2}$. Véase
        que $g = f - f(0)$ cumple
        $g\bigl(\frac x2\bigr) = \frac{g(x)}{2}$, dedúzcase que $g$ es
        aditiva y conclúyase $f(x) = cx + d$.
  \item Supóngase ahora solo la \emph{desigualdad}: $f$ \omterm{def:b1:continuity:continuous}{continua}
        con
        \[
        f\Bigl(\frac{x + y}{2}\Bigr) \leq \frac{f(x) +
        f(y)}{2} \qquad (x, y \in \R).
        \]
        Demuéstrese por inducción sobre $n$ que, para todos los pesos
        diádicos $\lambda = \frac{k}{2^n} \in \intcc{0}{1}$:
        \[
        f\bigl(\lambda x + (1 - \lambda)y\bigr) \leq \lambda
        f(x) + (1 - \lambda) f(y) .
        \]
  \item Extiéndase por \omterm{def:b1:continuity:continuous}{continuidad} y por la \omterm{def:b1:topology:dense}{densidad} de los
        diádicos (\cref{exo:b1:reals:8}) a todo
        $\lambda \in \intcc{0}{1}$: una función \omterm{def:b1:continuity:continuous}{continua} y convexa en
        el punto medio cumple la desigualdad de convexidad completa
        (la noción que se estudia sistemáticamente en
        el~\cref{ch:b1:derivative}).
  \item Véase que no se puede prescindir de la \omterm{def:b1:continuity:continuous}{continuidad}: una
        $f$ aditiva no lineal cumple la \emph{igualdad en el punto
        medio} de la pregunta 19 y, aun así, ninguna desigualdad de
        convexidad en ningún \omterm{prop:b1:reals:intervals}{intervalo} (pregunta 12). Moraleja: la
        convexidad en el punto medio es una propiedad de etapas
        numerables (los diádicos) y la convexidad, una del \omterm{def:b1:continuity:continuous}{continuo};
        la \omterm{def:b1:continuity:continuous}{continuidad} es el puente --- exactamente como en las
        partes I--II.
\end{enumerate}

\medskip
\noindent\textbf{Parte VI --- Últimas variaciones y síntesis.}
\begin{enumerate}[resume]
  \item Hállense todas las $f \colon \R \to \R$ \omterm{def:b1:continuity:continuous}{continuas} con
        $f(x + y) = f(x) + f(y) + xy$ \emph{(réstese la solución
        particular $\frac{x^2}{2}$)}.
  \item Demuéstrese: si $f \colon \R \to \R$ es \omterm{def:b1:continuity:continuous}{continua} y aditiva
        solo sobre un \emph{\omterm{def:b1:structures:subgroup}{subgrupo} \omterm{def:b1:topology:dense}{denso}} $G$ (es decir,
        $f(g + g') = f(g) + f(g')$ para $g, g' \in G$), entonces $f$
        es aditiva en $\R$. Más en general, dos funciones \omterm{def:b1:continuity:continuous}{continuas}
        que coinciden en un subconjunto \omterm{def:b1:topology:dense}{denso} de $\R$ son iguales.
  \item Síntesis, una frase para cada punto: (i) enúnciese de memoria
        la escalera de regularidad de la pregunta 8; (ii) explíquese
        por qué «gráfica \omterm{def:b1:topology:dense}{densa} en el plano» es la imagen mental
        correcta del fracaso de la regularidad; (iii) enumérense las
        cinco funciones clásicas caracterizadas en las partes IV y V y
        el único método que las atrapó a todas; (iv) nómbrense los dos
        lugares donde la \omterm{def:b1:topology:dense}{densidad} de $\Q$ (o de los diádicos) en $\R$
        sostuvo el argumento, y el lugar donde no pudo hacerlo
        (pregunta 13).
\end{enumerate}
\end{problem}
